\documentclass[10pt,a4paper]{amsart}
\usepackage{amsmath,amssymb,amsthm,graphicx}
\usepackage{verbatim}
\newenvironment{reference}{\comment}{\endcomment}
%\newenvironment{reference}{}{}
\newenvironment{slogan}{\comment}{\endcomment}
\newenvironment{history}{\comment}{\endcomment}

% For commutative diagrams we use Xy-pic
\usepackage[all]{xy}

% We use 2cell for 2-commutative diagrams.
\xyoption{2cell}
% Theorem environments.
%
\theoremstyle{plain}
\newtheorem{theorem}[subsection]{Theorem}
\newtheorem{proposition}[subsection]{Proposition}
\newtheorem{lemma}[subsection]{Lemma}

\theoremstyle{definition}
\newtheorem{definition}[subsection]{Definition}
\newtheorem{example}[subsection]{Example}
\newtheorem{exercise}[subsection]{Exercise}
\newtheorem{situation}[subsection]{Situation}

\theoremstyle{remark}
\newtheorem{remark}[subsection]{Remark}
\newtheorem{remarks}[subsection]{Remarks}

\numberwithin{equation}{subsection}

% Macros
%
\def\lim{\mathop{\mathrm{lim}}\nolimits}
\def\colim{\mathop{\mathrm{colim}}\nolimits}
\def\Spec{\mathop{\mathrm{Spec}}}
\def\Hom{\mathop{\mathrm{Hom}}\nolimits}
\def\Ext{\mathop{\mathrm{Ext}}\nolimits}
\def\SheafHom{\mathop{\mathcal{H}\!\mathit{om}}\nolimits}
\def\SheafExt{\mathop{\mathcal{E}\!\mathit{xt}}\nolimits}
\def\Sch{\mathit{Sch}}
\def\Mor{\mathop{\mathrm{Mor}}\nolimits}
\def\Ob{\mathop{\mathrm{Ob}}\nolimits}
\def\Sh{\mathop{\mathit{Sh}}\nolimits}
\def\NL{\mathop{N\!L}\nolimits}
\def\CH{\mathop{\mathrm{CH}}\nolimits}
\def\proetale{{pro\text{-}\acute{e}tale}}
\def\etale{{\acute{e}tale}}
\def\QCoh{\mathit{QCoh}}
\def\Ker{\mathop{\mathrm{Ker}}}
\def\Im{\mathop{\mathrm{Im}}}
\def\Coker{\mathop{\mathrm{Coker}}}
\def\Coim{\mathop{\mathrm{Coim}}}

% Boxtimes
%
\DeclareMathSymbol{\boxtimes}{\mathbin}{AMSa}{"02}

%
% Macros for moduli stacks/spaces
%
\def\QCohstack{\mathcal{QC}\!\mathit{oh}}
\def\Cohstack{\mathcal{C}\!\mathit{oh}}
\def\Spacesstack{\mathcal{S}\!\mathit{paces}}
\def\Quotfunctor{\mathrm{Quot}}
\def\Hilbfunctor{\mathrm{Hilb}}
\def\Curvesstack{\mathcal{C}\!\mathit{urves}}
\def\Polarizedstack{\mathcal{P}\!\mathit{olarized}}
\def\Complexesstack{\mathcal{C}\!\mathit{omplexes}}
% \Pic is the operator that assigns to X its picard group, usage \Pic(X)
% \Picardstack_{X/B} denotes the Picard stack of X over B
% \Picardfunctor_{X/B} denotes the Picard functor of X over B
\def\Pic{\mathop{\mathrm{Pic}}\nolimits}
\def\Picardstack{\mathcal{P}\!\mathit{ic}}
\def\Picardfunctor{\mathrm{Pic}}
\def\Deformationcategory{\mathcal{D}\!\mathit{ef}}
\usepackage{cleveref}
\crefname{originallemma}{original lemma}{original lemmas}
\Crefname{originallemma}{Original lemma}{Original lemmas}
\makeatletter
\newlabel{coherent-lemma-quasi-coherent-affine-cohomology-zero}{{01XB}{source}{}{sourcepointer-1}{}}
\newlabel{coherent-lemma-quasi-coherent-affine-cohomology-zero@cref}{{[originallemma][0][]01XB}{[1][0][]source}{}{}{}}
\makeatother
\crefname{originaldefinition}{original definition}{original definitions}
\Crefname{originaldefinition}{Original definition}{Original definitions}
\makeatletter
\newlabel{descent-definition-structure-sheaf}{{03DU}{source}{}{sourcepointer-2}{}}
\newlabel{descent-definition-structure-sheaf@cref}{{[originaldefinition][0][]03DU}{[1][0][]source}{}{}{}}
\makeatother
\crefname{originallemma}{original lemma}{original lemmas}
\Crefname{originallemma}{Original lemma}{Original lemmas}
\makeatletter
\newlabel{descent-lemma-compare-sites}{{070S}{source}{}{sourcepointer-3}{}}
\newlabel{descent-lemma-compare-sites@cref}{{[originallemma][0][]070S}{[1][0][]source}{}{}{}}
\makeatother
\crefname{originallemma}{original lemma}{original lemmas}
\Crefname{originallemma}{Original lemma}{Original lemmas}
\makeatletter
\newlabel{descent-lemma-sheaf-condition-holds}{{03DT}{source}{}{sourcepointer-4}{}}
\newlabel{descent-lemma-sheaf-condition-holds@cref}{{[originallemma][0][]03DT}{[1][0][]source}{}{}{}}
\makeatother
\crefname{originalproposition}{original proposition}{original propositions}
\Crefname{originalproposition}{Original proposition}{Original propositions}
\makeatletter
\newlabel{descent-proposition-equivalence-quasi-coherent}{{03DX}{source}{}{sourcepointer-5}{}}
\newlabel{descent-proposition-equivalence-quasi-coherent@cref}{{[originalproposition][0][]03DX}{[1][0][]source}{}{}{}}
\makeatother
\crefname{originalproposition}{original proposition}{original propositions}
\Crefname{originalproposition}{Original proposition}{Original propositions}
\makeatletter
\newlabel{descent-proposition-same-cohomology-quasi-coherent}{{03DW}{source}{}{sourcepointer-6}{}}
\newlabel{descent-proposition-same-cohomology-quasi-coherent@cref}{{[originalproposition][0][]03DW}{[1][0][]source}{}{}{}}
\makeatother
\crefname{originallemma}{original lemma}{original lemmas}
\Crefname{originallemma}{Original lemma}{Original lemmas}
\makeatletter
\newlabel{sites-cohomology-lemma-cohomology-over-U-trivial}{{0BKZ}{source}{}{sourcepointer-7}{}}
\newlabel{sites-cohomology-lemma-cohomology-over-U-trivial@cref}{{[originallemma][0][]0BKZ}{[1][0][]source}{}{}{}}
\makeatother
\title{Derived pullback to the small etale site is an equivalence on derived categories of quasi-coherent modules}
\author{The Stacks Project Authors}
\date{stacks-project:spaces-perfect.tex:local-snapshot-20261012; GNU FDL 1.2 or later}
\begin{document}
\maketitle
\noindent\textit{Editorial note.} All delivered mathematical text is the author's own, copied from the pinned Stacks Project chapter file spaces-perfect.tex; no statement or proof step was rewritten, reordered or removed.\par\medskip
\noindent\textit{Editorial note.} The counted proof is the author's complete argument: a module-level unit and counit check followed by an affine-cohomology computation of the derived pushforward and sheafification, closing both adjunction maps inside the delivered span\par\medskip
\noindent\textit{Editorial note.} The label lemma-derived-quasi-coherent-small-etale-site names the counted statement, and the proof environment immediately following it in the source is delivered verbatim as its proof.\par\medskip
\noindent\textit{Editorial note.} The counted statement is carried by the author's own lemma environment in the chapter source.\par\medskip
\noindent\textit{Allowed standard prerequisites (editorial).}\begin{itemize}
\item The comparison of the small Zariski and etale sites and the sheaf property of the associated etale structure sheaf from Descent
\item The equivalence of quasi-coherent modules on the small etale site with quasi-coherent modules on X, from Descent
\item Vanishing of higher cohomology of quasi-coherent sheaves on affine objects, and the site-level lemma expressing cohomology over U through the top cohomology sheaf
\end{itemize}\medskip
\ifdefined\hypertarget\hypertarget{sourcepointer-1}{}\fi
\noindent\textit{Source reference.} Original lemma 01XB: sources/stacks-chow/tags.txt, source lines 2512--2512. Prerequisite pointer invoked at proof line(s) 442; it names a standard fact of a sibling chapter that the counted argument uses, and every step of the counted proof built on it is retained verbatim in the delivered body.\par\medskip
\ifdefined\hypertarget\hypertarget{sourcepointer-2}{}\fi
\noindent\textit{Source reference.} Original definition 03DU: sources/stacks-chow/tags.txt, source lines 4433--4433. Prerequisite pointer invoked at proof line(s) 421; it names a standard fact of a sibling chapter that the counted argument uses, and every step of the counted proof built on it is retained verbatim in the delivered body.\par\medskip
\ifdefined\hypertarget\hypertarget{sourcepointer-3}{}\fi
\noindent\textit{Source reference.} Original lemma 070S: sources/stacks-chow/tags.txt, source lines 8966--8966. Prerequisite pointer invoked at proof line(s) 415; it names a standard fact of a sibling chapter that the counted argument uses, and every step of the counted proof built on it is retained verbatim in the delivered body.\par\medskip
\ifdefined\hypertarget\hypertarget{sourcepointer-4}{}\fi
\noindent\textit{Source reference.} Original lemma 03DT: sources/stacks-chow/tags.txt, source lines 4432--4432. Prerequisite pointer invoked at proof line(s) 423; it names a standard fact of a sibling chapter that the counted argument uses, and every step of the counted proof built on it is retained verbatim in the delivered body.\par\medskip
\ifdefined\hypertarget\hypertarget{sourcepointer-5}{}\fi
\noindent\textit{Source reference.} Original proposition 03DX: sources/stacks-chow/tags.txt, source lines 4436--4436. Prerequisite pointer invoked at proof line(s) 427; it names a standard fact of a sibling chapter that the counted argument uses, and every step of the counted proof built on it is retained verbatim in the delivered body.\par\medskip
\ifdefined\hypertarget\hypertarget{sourcepointer-6}{}\fi
\noindent\textit{Source reference.} Original proposition 03DW: sources/stacks-chow/tags.txt, source lines 4435--4435. Prerequisite pointer invoked at proof line(s) 439; it names a standard fact of a sibling chapter that the counted argument uses, and every step of the counted proof built on it is retained verbatim in the delivered body.\par\medskip
\ifdefined\hypertarget\hypertarget{sourcepointer-7}{}\fi
\noindent\textit{Source reference.} Original lemma 0BKZ: sources/stacks-chow/tags.txt, source lines 14577--14577. Prerequisite pointer invoked at proof line(s) 444; it names a standard fact of a sibling chapter that the counted argument uses, and every step of the counted proof built on it is retained verbatim in the delivered body.\par\medskip
\begin{lemma}
\label{lemma-derived-quasi-coherent-small-etale-site}
Let $X$ be a scheme. The functor
$\epsilon^* : D_\QCoh(\mathcal{O}_X) \to
D_\QCoh(\mathcal{O}_\etale)$
defined above is an equivalence.
\end{lemma}

\begin{proof}
We will prove this by showing the functor
$R\epsilon_* : D(\mathcal{O}_\etale) \to D(\mathcal{O}_X)$
induces a quasi-inverse. We will use freely that $\epsilon_*$
is given by restriction to $X_{Zar} \subset X_\etale$ and the description of
$\epsilon^* = \text{id}_{small, \etale, Zar}^*$
in Descent, Lemma \ref{descent-lemma-compare-sites}.

\medskip\noindent
For a quasi-coherent $\mathcal{O}_X$-module $\mathcal{F}$ the adjunction map
$\mathcal{F} \to \epsilon_*\epsilon^*\mathcal{F}$ is an isomorphism by
the fact that $\mathcal{F}^a$
(Descent, Definition \ref{descent-definition-structure-sheaf})
is a sheaf as proved in
Descent, Lemma \ref{descent-lemma-sheaf-condition-holds}.
Conversely, every quasi-coherent $\mathcal{O}_\etale$-module
$\mathcal{H}$ is of the form $\epsilon^*\mathcal{F}$ for some quasi-coherent
$\mathcal{O}_X$-module $\mathcal{F}$, see
Descent, Proposition \ref{descent-proposition-equivalence-quasi-coherent}.
Then $\mathcal{F} = \epsilon_*\mathcal{H}$ by what we just said and
we conclude that the adjunction map
$\epsilon^*\epsilon_*\mathcal{H} \to \mathcal{H}$ is an isomorphism for all
quasi-coherent $\mathcal{O}_\etale$-modules $\mathcal{H}$.

\medskip\noindent
Let $E$ be an object of $D_\QCoh(\mathcal{O}_\etale)$
and denote $\mathcal{H}^q = H^q(E)$ its $q$th cohomology
sheaf. Let $\mathcal{B}$ be the set of affine objects of $X_\etale$.
Then $H^p(U, \mathcal{H}^q) = 0$ for all $p > 0$, all $q \in \mathbf{Z}$,
and all $U \in \mathcal{B}$, see
Descent, Proposition \ref{descent-proposition-same-cohomology-quasi-coherent}
and
Cohomology of Schemes, Lemma
\ref{coherent-lemma-quasi-coherent-affine-cohomology-zero}.
By Cohomology on Sites, Lemma
\ref{sites-cohomology-lemma-cohomology-over-U-trivial}
this means that
$$
H^q(U, E) = H^0(U, \mathcal{H}^q)
$$
for all $U \in \mathcal{B}$. In particular, we find that this holds
for affine opens $U \subset X$. It follows that the $q$th cohomology of
$R\epsilon_*E$ over $U$ is the value of the sheaf $\epsilon_*\mathcal{H}^q$
over $U$. Applying sheafification we obtain
$$
H^q(R\epsilon_*E) = \epsilon_*\mathcal{H}^q
$$
which in particular shows that $R\epsilon_*$ induces a functor
$D_\QCoh(\mathcal{O}_\etale) \to D_\QCoh(\mathcal{O}_X)$.
Since $\epsilon^*$ is exact we then obtain
$H^q(\epsilon^*R\epsilon_*E) = \epsilon^*\epsilon_*\mathcal{H}^q =
\mathcal{H}^q$ (by discussion above). Thus the adjunction map
$\epsilon^*R\epsilon_*E \to E$ is an isomorphism.

\medskip\noindent
Conversely, for $F \in D_\QCoh(\mathcal{O}_X)$ the
adjunction map $F \to R\epsilon_*\epsilon^*F$
is an isomorphism for the same reason, i.e., because
the cohomology sheaves of $R\epsilon_*\epsilon^*F$
are isomorphic to
$\epsilon_*H^m(\epsilon^*F) = \epsilon_*\epsilon^*H^m(F) = H^m(F)$.
\end{proof}

\end{document}
